% res_rejoin.tex -- DRAFT 23 Oct 2026: a departed carrier rejoins its own plate (rejoin_departed; twin only).
% Source: card docs/experiments/2026-10-23_rejoin_departed.md (critic 22 Oct answered; E0 bars frozen 23 Oct 00:08, E1
% 00:42, E2-E4 05:13, E5 07:57; critic on the 23 Oct fixes 06:25 "run it", its M2 fixed and M1 logged). Registry: E0 R2445;
% E1 R2457/R2458 (FAIL, 1.5 s blend), R2466/R2467 (VOID as E1b, FAIL as E1), R2578/R2579 (blend fix, PASS), R2592
% (05:59 fixes, PASS), R2596 (hold fix = final tree, PASS); E2 R2584 (FAIL 10.61 deg), R2589 (05:59, PASS), R2594/R2595
% (final, PASS); E3 R2585 (FAIL 1.179 m), R2590 (05:59, PASS); E4 R2586 (FAIL 1.709 m), R2591 (05:59, PASS).
% Code trees: (i) be34f95 + rejoin_blend_s (the rejoin's own 10 s blend); (ii) (i) + the 05:59 fixes (_rejoin_cap,
% REJOIN_OUT_MARGIN_M and the exact way out, the staging bound STAGE_TRACK_M); (iii) (ii) + the path held for
% rejoin_blend_s and the cap's binding logged (the final tree). E3 and E4 flew on (ii): the hold fix acts only through a
% fold-in's blend, which neither reaches (E3 aborts in the descent, E4's ATTACH is refused) -- an inference from the code,
% not a re-fly. E5 (the pair leaning ~3 deg toward / away from the rejoin plate, world three_rigid_ground_rig_T_push_L100,
% Wesley 23 Oct (10); bars frozen 07:57, as E2's): toward R2602 (+3.26 deg at the weld, PASS 3.12 deg) and R2712 (+2.95,
% FAIL 5.28 deg, over 5 from weld +0.22 to +0.41 s: the weld-height cap dipped the rejoiner's reference 1.1 cm, 2.044 ->
% 2.033 m, the critic's M1); away R2603 (-2.67, PASS 2.91) and R2713 (-3.11, PASS 3.37; lean before LAND 1.88 / 1.93
% against 2, the torque on to LAND). Wesley 23 Oct ~14:05: E5 reported as is, 1 of 2 clean toward the plate; no change
% to the cap. Every other E5 bar passes in all four runs.
% Must not say: anything for the rig (the twin's rods are rigid, can push, have no collision geometry, and its joint welds
% at any distance; the rig's magnet catches only on contact, and on the rig the collaborator's controller flies every
% attach). The cap patches a twin mismatch (on a T the rider's rigid rod height IS the ring's pitch), during the blend only.

\subsection{A departed carrier rejoins}
\label{sec:res_rejoin}

This section asks whether a carrier that has left the ring can be flown back to its own plate, re-attached and folded
back into the team, so that the team returns to its size with the ring level. It is the step that closes the chains of
Section~\ref{sec:res_scale}, and it is a twin result only: the twin's joint welds the rod wherever its end is, while
the rig's magnet catches only on contact, and on the rig every attach is flown by the collaborator's controller
(Section~\ref{sec:res_join}). The bars were frozen before each check flew.

\paragraph{How the rejoin flies.} % [Method: move to Section~\ref{sec:reconfig} when this section is placed.]
A freed carrier is not landed. It holds a staging point on its own plate's bearing, 0.3\,m beyond the formation's
reach, which moves and turns with the planned ring and stays inside the position limit. On the operator's ATTACH,
accepted only while the path is held, it flies the newcomer's approach (a climb, a transit over its plate and a
descent at the attach speed), aimed so that its rod's end lands on the plate point it left, with a loop on the
measured end's height. When the end has stayed within 1.5\,cm of the plate point for 0.3\,s and moves slower than
2\,cm/s, the planner asks the twin's joint to weld, waits for the joint to report it, and only then grows its
problem by one UAV, with the remembered plate point and the rod re-measured from the pivot to that point. The fold-in
is the newcomer's: the rejoiner from a vertical rod and a near-zero pull, blended to its share over 10\,s, with the
path held through the blend and the rejoiner's reference held at or below its weld height. A LAND or a drop during a
rejoin sends the rejoiner up, out past the clear radius and down alone.

Table~\ref{tab:rejoin} gives the fold-ins, Table~\ref{tab:rejoin_neighbours} the checks around them. A rejoin is
scored in stages: the rod's end on its plate within 30\,s of the rejoin's start; the joint attached within 0.5\,s of
that; and the fold-in, with the ring's tilt at most 5$^\circ$ from the weld to 15\,s after it, the end within 2\,cm of
its plate through the blend, the rejoiner's planned pull at the blend's end within 10\,\% of its balanced share, the
ring's lean over the last 10\,s before LAND at most 2$^\circ$, and every UAV landed with no drop or abort. Every
neighbouring check passed, and on the final code tree every fold-in but one. The rejoin onto the square (four UAVs back from
three) and onto the T (three back from the opposite pair) both brought the rod's end to its plate inside the 30\,s
allowed, the joint answered within a few hundredths of a second, and the ring stayed well inside its tilt bar through
the fold-in. With the T's pair leaning toward the rejoin plate, one of two fold-ins tipped the ring just past the bar,
briefly; leaning away, both passed, with the ring's lean before LAND the closest to its bar.

\begin{table}[htbp]
\centering
\footnotesize
\setlength{\tabcolsep}{2.4pt}
\caption[A departed carrier rejoins]{A departed carrier rejoins its own plate in the rig twin: the 0.86\,kg ring at
1.0\,m on 1\,m cables, hover, against the bars frozen before the runs (in brackets). Square: the UAV on plate 9 leaves
the square by a ramped release and rejoins it. T: the T's rider (plate 3) is released unannounced, the opposite pair
carries on, and the rider rejoins. Code trees, in order: (i) the rejoin's own 10\,s blend; (ii) with the reference
held at or below its weld height through the blend, the exact way out and the staging bound; (iii) the final tree, with
the path held through the whole blend. Leaning: a steady torque on the ring leans the T's pair about 3$^\circ$ toward
the rejoin plate or away from it. Ranges over $n$ runs. Share: the rejoiner's planned pull at the blend's end against
its balanced share. Lean: the ring's mean tilt over the last 10\,s before LAND (when leaning, with the push still on).
$^a$\,The blend ran for 1.5\,s, not 10\,s (a configuration error, fixed by tree (i)); two more runs with the same short
blend tilted 5.75--5.82$^\circ$. $^b$\,The rejoiner's reference rose 4.6\,cm as the blend swung its rod out, and its
rigid rod lifted the T's free side. $^c$\,5.28$^\circ$ in one run, over the bar for 0.19\,s: at the fold-in the weld-height hold
dipped the rejoiner's reference 1.1\,cm while the pair leaned toward the plate.}
\label{tab:rejoin}
\begin{tabular}{llccccccc}
\toprule
 & & & On its plate & Joint & Ring tilt, weld & End from & & Lean \\
Rejoin & Tree & $n$ & after start (s) & after (s) & to $+$15\,s ($^\circ$) & plate (cm) & Share & ($^\circ$) \\
 & & & [30] & [0.5] & [5] & [2] & [10\,\%] & [2] \\
\midrule
Square & (i) & 2 & 24.86--24.87 & 0.029--0.033 & 1.27--1.30 & 1.3--1.4 & $+$2.5\,\% & 0.08--0.09 \\
 & (ii) & 1 & 24.81 & 0.036 & 1.45 & 1.4 & $+$2.5\,\% & 0.09 \\
 & (iii) & 1 & 23.87 & 0.030 & 1.55 & 1.6 & $+$2.1\,\% & 0.09 \\
T & (ii) & 1 & 17.47 & 0.021 & 3.02 & 1.4 & $-$2\,\% & 0.71 \\
 & (iii) & 2 & 16.58--17.66 & 0.012--0.015 & 1.45--2.25 & 1.3--1.5 & $-$1 to 0\,\% & 0.72--0.79 \\
T, leaning toward & (iii) & 2 & 15.56--16.47 & 0.015--0.018 & 3.12--5.28$^c$ & 1.3--1.5 & $-$1 to 0\,\% & 0.49--0.52 \\
T, leaning away & (iii) & 2 & 17.67--18.56 & 0.015--0.018 & 2.91--3.37 & 1.4--1.5 & $-$1 to 0\,\% & 1.88--1.93 \\
\midrule
\multicolumn{9}{l}{\emph{Before the fixes (failed, kept as results)}} \\
Square$^a$ & -- & 2 & 24.87--25.80 & 0.024--0.027 & 5.32--5.83 & 1.3--1.4 & $-$2\,\% & 0.08--0.09 \\
T & (i) & 1 & 17.47 & 0.021 & 10.61$^b$ & 1.3 & 0\,\% & 0.74 \\
\bottomrule
\end{tabular}
\end{table}
% Square (i) R2578/R2579; (ii) R2592; (iii) R2596. T (ii) R2589; (iii) R2594/R2595. Leaning toward: R2602 / R2712 (lab
% batches 16 / 17; pair lean at the weld +3.26 / +2.95 deg; share +0 / -0.9 %); away: R2603 / R2713 (-2.67 / -3.11 deg;
% -1 / +0.1 %). The cap bound 101 ticks in R2602/R2603, the plan 2.7 / 2.3 cm above it at the last.
% Short blend: R2457/R2458 (scored E1, FAIL); footnote a's two runs R2466/R2467 (VOID as E1b, FAIL as E1 repeats,
% 5.82/5.75 deg).
% T before the fixes: R2584 (tree (i)). Lean column: the registry rows (R2457/R2458 0.09/0.08; R2466/R2467 not recorded
% beyond '<= 2'). Every run landed every UAV, no drop, abort or refusal.

The two failures before the fixes are results. With a 1.5\,s blend the rejoiner took its share almost at once and the
ring tipped just past the bar for a fraction of a second; the full 10\,s blend removed it. On the T the failure was a
twin mismatch: on a T only the rider holds the ring's pitch about the pair's line, and in the twin the rider's rigid
rod makes its height the ring's pitch. As the blend swung the rejoiner's rod out from vertical, its planned height
rose and its rod lifted the T's free side while it carried almost nothing. Holding its reference at or below its weld
height through the blend removed the lift. The hold patches the twin's rigid rod during the blend only; on the rig's
slack tether the mismatch does not arise. Its own small dip at the fold-in, and its release at the blend's end, give
the T's remaining tilt. The hold clamps the reference outside the planner's problem, so its dip is a disturbance of its
own, and it is least sure of itself when the pair leans toward the rejoin plate. Flown with the pair leaning about
3$^\circ$ either way, the fold-ins stayed inside the bar but for one run leaning toward the plate, where the hold's dip
at the fold-in tipped the ring just past it, briefly. That run is reported as it flew; the hold was not changed. With
the pair leaning away, the push that leans it, left on to the end, keeps the ring's lean before LAND close to its bar.

\begin{table}[htbp]
\centering
\footnotesize
\caption[Around the rejoin]{The checks around the rejoin in the rig twin, one run each, against the bars frozen
before the runs. $^a$\,On the floor before the flight the twin's joint welded the rod with its end where it lay, off
the plate; in flight the weld gate prevents this.}
\label{tab:rejoin_neighbours}
\begin{tabularx}{\textwidth}{>{\raggedright\arraybackslash}p{3.4cm}>{\raggedright\arraybackslash}X>{\raggedright\arraybackslash}X>{\raggedright\arraybackslash}p{1.3cm}}
\toprule
Check & Bar & Measured & Verdict \\
\midrule
A rod re-welded on the floor, then the square lifts and hovers (first tree)
  & joint released and re-welded; hover with ring tilt at most 2$^\circ$; every UAV lands
  & both joint states within 0.6\,s; ring tilt 5.6$^\circ$ (median), 6.7$^\circ$ (peak)$^a$; all landed & fails \\
LAND 1.5\,s into the rejoin's descent (tree ii)
  & the rejoin aborts; it climbs until its end is 0.25\,m over its plate before moving out; it lands alone at least
  1.2\,m from the ring centre; no UAV in the air when the fleet is declared landed
  & aborted 0.10\,s after LAND; end 0.27\,m over its plate; landed 1.245\,m out; within 1\,cm of its rest at the
  declaration & passes \\
Staging through a figure-eight lap of the T, with an ATTACH mid-lap (tree ii)
  & the ATTACH refused; the staged UAV inside $\pm$1.6\,m, at least 0.6\,m from the ring centre and 0.8\,m from every
  carrier
  & refused; at most 1.575\,m out along $x$; 0.739\,m from the centre, 1.028\,m from the nearest carrier & passes \\
\midrule
\multicolumn{4}{l}{\emph{Before the fixes (first tree)}} \\
LAND in the descent & as above & landed 1.179\,m out: the way out stopped short & failed \\
Staging on a path & as above & reached 1.709\,m: the staging reference overshot its bound & failed \\
\bottomrule
\end{tabularx}
\end{table}
% E0 R2445 (JOINT_DETACHED 0.004 s after RELEASE, JOINT_ATTACHED 0.58 s after the publish; the re-welded quad sat
% 0.761 m out against 0.733 for its neighbours). E3 R2590 (R2585 before). E4 R2591 (R2586 before); the T's path error in
% the lap 45.5 mm, ring tilt peak 2.74 deg (reported, no bar).

The rejoin also closes the eight-UAV chain of Section~\ref{sec:res_scale}: on the final tree four freed UAVs
rejoined a 0.5\,m ring one at a time, each welded within 20\,s of its start, with every rod's end within 2\,cm of its
plate through its blend (exploratory, one run).

\pendingfig{fig:rejoin}{A freed rider rejoins the T in the rig twin (final code tree, hover, 1\,m cables). (a) The
distance of its rod's end from the plate point it left, against the 1.5\,cm weld gate, with the start, the end on the
plate, the weld and the end of the blend marked; (b) the rejoiner's height reference through the 10\,s blend, with
the cap at its weld height; (c) the ring's tilt against the 5$^\circ$ bar; (d) the rejoiner's planned pull rising to
its share.}{a final-tree rejoin onto the T}
% Data: R2594 (or R2595): run.csv (tether tip poses, payload tilt), the planner log (the rejoiner's reference z and
% planned tension), events.csv (REJOIN_START, REJOIN_ON_PLATE, JOINT_ATTACHED, REJOIN_WELD, REJOIN_BLEND_OVER). Tool: a
% new thesis_sim_figures function through the thesis-figure skill; R2596 (the square) as a second column if it earns it.

\paragraph{Limits.} A rejoin flies only while the path is held; on a path the staged UAV follows its plate and the
ATTACH is refused. The weld waits for the released rod's swing to die by itself: the loop on the measured end trims its
height, not its swing. And the twin flatters the rejoin: its rods are rigid, can push and have no collision geometry,
and its joint welds within the gate wherever the end is.

\paragraph{Answer.} In the twin, a carrier that has left the ring flies back to its own plate, is re-welded and is
folded back in with the ring level, onto a square and onto a T, and the neighbouring checks of a LAND mid-rejoin and of
staging on a path pass. It needed two fixes that the twin's rigid rods and one configuration error forced. With the
T's pair leaning toward the rejoin plate, one fold-in in two went just past the tilt bar, from the dip of the hold that
the rigid rods forced. It is not a rig result.
